\documentclass{article}

\usepackage[margin=1in]{geometry}
\usepackage[utf8]{inputenc}
\usepackage[T1]{fontenc}
\usepackage{lmodern}
\usepackage{booktabs}
\usepackage{array}
\usepackage{hyperref}
\usepackage{url}
\usepackage{microtype}
\setlength{\emergencystretch}{3em}
\hypersetup{
  pdftitle={Supplement: Offensive CTFs and SOC Investigations Stress Different Agent Behaviors},
  pdfauthor={Paul Kassianik},
  pdfsubject={Cybersecurity benchmark evaluation supplement},
  pdfkeywords={cybersecurity benchmarks, AI agents, Cybench, BOTS, Inspect, SOC}
}

\title{Supplement: Offensive CTFs and SOC Investigations Stress Different Agent Behaviors}
\author{Paul Kassianik\\Stealth Company}
\date{Short Paper Submission\\\today}

\begin{document}
\maketitle

\section{Scope}

This supplement records aggregate provenance and robustness checks for the short paper.  It is intended for venues that allow appendices or supplemental PDFs.  It does not include raw Inspect transcripts, private logs, credentials, source code, or benchmark artifacts.  Logical source paths below refer to \path{source-eval-repo/...}; local machine paths are intentionally omitted.

No additional benchmark, Docker, Kubernetes, or Inspect workload was launched for this supplement.  The analyses below reuse existing successful Cybench and BOTSv1 Inspect logs and generated chart artifacts.

\section{Source provenance}

The quantitative tables in the paper are derived from these sanitized logical sources:
\begin{itemize}
  \item \path{source-eval-repo/charts/cybench_cost_table.csv}
  \item \path{source-eval-repo/charts/cybench_report.html}
  \item \path{source-eval-repo/charts/botsv1_report.html}
  \item raw \path{source-eval-repo/logs/*.eval} files listed in \path{tables/provenance.md}
\end{itemize}
At audit time the source evaluation repository HEAD was \texttt{0cf3bd8}; generated chart/report artifacts had local working-tree updates and should be checksummed or committed before any public artifact release.

\section{Uncertainty method}

The paper reports descriptive robustness intervals rather than formal population-level significance tests.  For Cybench, each bootstrap replicate resamples challenge IDs with replacement and retains the three epochs for each selected challenge.  For BOTSv1, each replicate resamples question IDs with replacement and retains the three epochs and official point weights for each selected question.  Paired contrasts resample only task IDs common to both compared logs.

\section{Cybench robustness}

\begin{table}[h]
\centering
\caption{Cybench challenge-level bootstrap intervals.}
\begin{tabular}{lrr}
\toprule
Model & pass@1 & 95\% bootstrap interval \\
\midrule
GPT-5.5 + autocompact & 94.0\% & 87.2--99.2\% \\
DeepSeek v4 Flash + autocompact (\$2.10 cap) & 86.3\% & 76.9--94.9\% \\
Claude Opus 4.7 + autocompact & 70.9\% & 59.0--82.9\% \\
DeepSeek v4 Flash + autocompact (\$0.80 cap) & 64.1\% & 50.4--76.9\% \\
Kimi K2.6 + autocompact & 55.6\% & 41.0--70.1\% \\
DeepSeek v4 Pro + autocompact & 43.9\% & 29.8--57.9\% \\
GPT-5.4 Mini + autocompact & 29.1\% & 16.2--42.7\% \\
\bottomrule
\end{tabular}
\end{table}

The most defensible Cybench scaling claim is the paired DeepSeek v4 Flash budget comparison: the \$2.10-cap run exceeds the \$0.80-cap run by 22.2 percentage points, with a challenge-level paired bootstrap interval of +12.8 to +32.5 percentage points.  GPT-5.5 exceeds the higher-budget DeepSeek v4 Flash row by 7.7 percentage points, with interval +1.7 to +15.4 percentage points.

\section{BOTSv1 robustness}

\begin{table}[h]
\centering
\caption{BOTSv1 question-level bootstrap intervals.}
\begin{tabular}{lrr}
\toprule
Model & BOTS points & 95\% bootstrap interval \\
\midrule
Claude Opus 4.8 & 93.9\% & 82.7--99.5\% \\
GPT-5.5 high effort & 81.4\% & 67.4--94.1\% \\
GPT-5.5 & 81.0\% & 64.1--96.4\% \\
DeepSeek v4 Pro & 77.8\% & 61.0--93.1\% \\
DeepSeek v4 Flash (\$4.20 cap) & 73.9\% & 53.2--91.9\% \\
DeepSeek v4 Flash & 73.0\% & 52.7--91.1\% \\
GPT-5.4 Mini & 27.5\% & 14.9--42.3\% \\
\bottomrule
\end{tabular}
\end{table}

Claude Opus 4.8's paired margin over the standard GPT-5.5 row is 12.8 percentage points, with a question-level paired bootstrap interval of +0.9 to +27.4 percentage points.  Its paired margin over GPT-5.5 high effort is 12.5 percentage points, with interval +2.4 to +23.7 percentage points.  The DeepSeek v4 Flash \$4.20-cap row exceeds the \$2.10-cap row by only 0.9 percentage points, with interval -4.9 to +6.9 percentage points, so the budget change should be described as inconclusive.

\section{BOTSv1 sequential-context check}

The BOTSv1 task excludes the official warm-up question and evaluates 31 scored questions.  Twenty-three of those questions have explicit dependencies on earlier case facts and receive prerequisite context in the prompt.  This matches the benchmark's sequential investigation structure, but it should be interpreted as follow-up investigation with known case state rather than cold-start reconstruction.

For Claude Opus 4.8, independent questions score 1,445/1,500 points (96.3\%) and dependent questions score 8,221.7/8,800 points (93.4\%).  The standard GPT-5.5 row scores 1,345/1,500 points (89.7\%) on independent questions and 7,000/8,800 points (79.5\%) on dependent questions.  Claude's mean-reduced point loss decomposes into 133.3 points from hint penalties and 500 points from misses; GPT-5.5 loses 205 points to hint penalties and 1,750 points to misses.  These aggregates support the paper's claim that SOC benchmark outcomes depend on evidence use and answer correctness, not simply on raw tool-call volume.

\section{Experiment-budget decision}

The no-cost robustness checks above should be preferred before spending the \$500 experiment budget.  If a venue reviewer specifically challenges prerequisite context, the highest-value follow-up would be a BOTSv1 \texttt{prereq\_context=false} ablation for the top model, run through Inspect under the approved non-prod Kubernetes namespace.  If a reviewer challenges the Claude version mismatch across task families, a Claude Opus 4.8 Cybench run would be the next candidate.  Neither experiment is necessary for the current short-paper claim, which is intentionally framed as an observational evaluation audit rather than a controlled model-only leaderboard.

\end{document}
